\section{A tour of the code, file by file}\label{app:code}

\begin{figure}[H]
\centering
\footnotesize
\begin{tabularx}{\linewidth}{@{}>{\raggedright\arraybackslash}p{7.6cm}X@{}}
\toprule
\textbf{Phase 1 (GPU, Python)} & \\
\file{roba/multipliers.py} & the eight multipliers in PyTorch, as per-number ``pieces''\\
\file{roba/reference.py} & the same eight with plain Python bit operations (independent check)\\
\file{roba/matmul.py} & approximate matrix multiply and convolution from a few exact ones\\
\file{roba/nn.py} & layers with switchable multiplier and number format; MLP, LeNet-5; calibration\\
\file{roba/data.py} & reads MNIST / Fashion-MNIST and normalises the pixels\\
\file{roba/cost.py} & the early cycle and energy \emph{estimate} (replaced by Ibex measurements)\\
\file{experiments/run.py} & stages errors, train, eval, finetune $\to$ \file{results/<dataset>/}\\
\midrule
\textbf{Shared C library} & \\
\file{csrc/roba.c}, \file{roba.h} & the multipliers in C, one function each (used by the ``call'' builds)\\
\file{csrc/roba\_core.h} & the same routines as inline functions, for the fused loop\\
\file{csrc/roba\_core\_test.c} & lets the Python tests call the inline routines\\
\midrule
\textbf{Phase 2 (Ibex)} & \\
\file{hw/setup.sh}, \file{build\_sim.sh}, \file{run\_sim.sh} & install tools, build the simulator, run a program\\
\file{hw/patches/0001-ibex-lod-instruction.patch} & the LOD instruction (Section~\ref{sec:lodpatch})\\
\file{hw/sw/lod\_test.c} & checks LOD against a software loop\\
\file{hw/sw/mult\_bench.c} & cycles per multiplication\\
\file{hw/sw/dnn.c} & the integer networks for the chip, every method\\
\file{hw/sw/util.h}, \file{link\_1mb.ld}, \file{Makefile} & printing, cycle counter, memory map, builds\\
\file{hw/run\_ibex.py} & the whole Phase 2 pipeline (Figure~\ref{fig:flow})\\
\file{hw/profile\_ibex.py} & where the cycles go (instruction trace of one picture)\\
\midrule
\textbf{Phase 3 (MLPerf Tiny)} & \\
\file{tiny/data.py} & CIFAR-10, Speech Commands (MFCC) and ToyADMOS (log-mel slices), with augmentation\\
\file{tiny/models.py} & ResNet-8, DS-CNN and the autoencoder, each with a layer list (\code{spec()})\\
\file{tiny/qnet.py} & int4/8/16 quantization, QAT model, export of the integer program, bit-exact reference\\
\file{tiny/run.py} & stages train, qat, eval for 3 seeds $\to$ \file{results/tiny/<dataset>/}\\
\file{hw/sw/tiny.c}, \file{tiny\_kernels.h} & the generic integer network for Ibex: convolution, shortcuts, pooling\\
\file{hw/run\_tiny.py} & writes \file{net.h} from the program, builds, simulates, checks every output\\
\file{tests/test\_tiny.py} & 33 tests of the integer layers, the square identity and input quantization\\
\midrule
\textbf{Checks and presentation} & \\
\file{tests/test\_multipliers.py} & 35 tests: PyTorch = Python = C for every multiplier\\
\file{experiments/build\_site.py}, \file{site/template.html} & build the slide deck \file{index.html}\\
\file{doc/make\_figures.py}, \file{doc/roba\_explained.tex} & this document\\
\bottomrule
\end{tabularx}
\caption{What each file does.}
\end{figure}

\subsection{\file{roba/multipliers.py}: one description for every RoBA variant}

Each RoBA variant is just a list of rounding rules, one per step:
\begin{lstlisting}[language=Python]
ROBA_ITERATIONS = {
    "roba_dn1":    [round_dn],
    "roba_updn":   [round_updn],
    "roba_dn2":    [round_dn, round_dn],
    "roba_updndn": [round_updn, round_dn],
    "roba_updn2":  [round_updn, round_updn],
}
\end{lstlisting}
\code{factors(x, mult, side)} turns a number into its pieces (Section~\ref{sec:factors}): for each step it
rounds the current leftover \code{r}, emits \code{[rr, r - rr]} for the left number or \code{[r, rr]} for the
right number, and continues with the leftover \code{r - rr}. Summing the products of matching pieces gives the
approximate product. \code{round\_dn} and \code{round\_updn} use \code{torch.frexp}, which splits a number into
$m\times2^e$; that is the GPU's way of finding the leading one.

\subsection{\file{roba/nn.py}: switchable layers}

\code{ALinear} and \code{AConv2d} behave like PyTorch's layers but have two settings:
\code{domain} (\code{None} = normal float32, \code{"float"}, \code{"trunc"}, \code{8}, \code{16}) and
\code{mult} (one of the eight). In \code{\_forward}, for int8 it quantizes inputs and weights
(\code{fake\_quant}), computes the exact product (which carries the gradient) and the approximate one
(computed without gradient), and returns \code{ref + (acc - ref).detach()}: the value of the approximate
result with the gradient of the exact one, the straight-through estimator. \code{calibrate()} records the
largest input of every layer on 5{,}000 training pictures, which sets the activation scales.

\subsection{\file{experiments/run.py}: the GPU stages}

\begin{lstlisting}[language=bash]
python3 experiments/run.py --dataset mnist   --stage all                              # ~7 min
python3 experiments/run.py --dataset fashion --stage all --epochs 30 --ft-epochs 5    # ~10 min
\end{lstlisting}
It writes \file{results/<dataset>/eval\_<model>.csv} (one row per number format and multiplier: accuracy,
loss, agreement with exact) and \file{finetune\_<model>.csv} (accuracy before and after fine-tuning), and saves
every checkpoint in \file{checkpoints/<dataset>/} (e.g.\ \file{mlp\_ft\_int8\_roba\_dn1.pt}), which Phase~2
reads.

\subsection{\file{hw/sw/dnn.c}: one program, many methods}

The same C file is compiled differently for each configuration. The \code{-D} flags pick the method:

\begin{table}[H]
\centering
\small
\begin{tabularx}{\linewidth}{lX}
\toprule
Flags & What the multiply-accumulate does\\
\midrule
\code{-DMULT\_NZ=roba\_dn1\_nz} & fused (default): \code{mac\_g}/\code{mac\_r} with the routine inlined from
\file{roba\_core.h}, sign-magnitude weight bytes (\code{WS0}, \dots, built in \code{tables\_init})\\
\code{-DMAC\_CALL -DMULT\_U=roba\_dn1\_u} & paper style: \code{mul\_s()} calls the routine in \file{roba.c} once per
product\\
\code{-DSQTAB} & table of squares: \code{STAB[509]} built in \code{tables\_init}, $\sum S(w)$ per neuron in
\code{model\_init}, $\sum S(x)$ per layer input\\
\code{-DZERO\_SKIP} & (with any of the above) gather the non-zero inputs (\code{gather\_nz}), or skip zero taps in
convolutions\\
\code{INPUT\_NONNEG} (in \file{model.h}) & raw pixels: the first layer also uses the ``input $\ge 0$'' MAC\\
\code{C0\_PER\_POSITION} (in \file{model.h}) & LeNet-5 with raw pixels: conv1 bias per output position\\
\bottomrule
\end{tabularx}
\caption{Build options of \file{hw/sw/dnn.c}.}
\end{table}

The network itself is written once, with a macro for a fully connected layer:
\begin{lstlisting}[language=C]
static void forward(const int8_t *x, uint32_t *lc)       /* the MLP */
{
    uint32_t t = (uint32_t)cycles(), u;
    FC(x, 784, 0, 256, h1, 0, MAC_IN);                   /* layer 0: 784 -> 256 */
    u = (uint32_t)cycles(); lc[0] = u - t; t = u;        /* cycles of layer 0   */
    FC(h1, 256, 1, 128, h2, 0, MAC_R);                   /* layer 1: 256 -> 128 */
    u = (uint32_t)cycles(); lc[1] = u - t; t = u;
    FC(h2, 128, 2, 10, z, 1, MAC_R);                     /* layer 2: 128 -> 10, keep logits */
    u = (uint32_t)cycles(); lc[2] = u - t;
}
\end{lstlisting}
\code{FC} expands to the right loop for the chosen method (plain, zero-skipping, table, or table with
skipping), followed by \code{rescale()} and \code{relu\_q()} for every output. \code{main()} builds the tables,
switches on the cycle counter, runs \code{forward} on every picture in \file{images.h}, and prints the
\code{IMG,\dots} line.

\subsection{\file{hw/run\_ibex.py}: the pipeline}

\begin{itemize}
  \item \code{int\_model()}: checkpoint $\to$ int8 weights, $M$, $C$ (and the raw-pixel fold).
  \item \code{quantize\_images()}: pictures $\to$ int8 (normalised or raw).
  \item \code{ref\_forward()}: the bit-exact Python copy of \file{dnn.c}; \code{logit\_hash()} the same
        fingerprint as the C code.
  \item \code{write\_model\_h()}, \code{write\_images\_h()}: the C headers.
  \item \code{compile\_elf()}: GCC + the ``no \code{mul}/\code{div}'' check. \code{run\_sim()}: one Verilator run.
  \item \code{prepare\_jobs()}, \code{run\_job()}, \code{stage\_dnn()}: build every configuration, run the
        batches in parallel (\code{--jobs 48}), reuse a log only if the program's SHA-256 matches, compare with
        the reference, write the CSV files.
\end{itemize}
A configuration name is a multiplier plus flags: \code{roba\_dn1}, \code{roba\_dn1+call}, \code{sqtab+zs+raw},
and so on.

\subsection{The CSV result files}

\begin{table}[H]
\centering
\small
\begin{tabularx}{\linewidth}{lX}
\toprule
File & Contents\\
\midrule
\file{results/<ds>/multiplier\_errors.csv} & error statistics per multiplier (8 and 16 bits)\\
\file{results/<ds>/eval\_<model>.csv} & GPU accuracy and loss per number format and multiplier\\
\file{results/<ds>/finetune\_<model>.csv} & GPU accuracy before/after fine-tuning\\
\file{results/ibex/mult\_bench.csv} & cycles and instructions per multiplication\\
\file{results/ibex/dnn\_summary.csv} & fused multipliers and the table: cycles per picture and per layer, accuracy,
match with the reference\\
\file{results/ibex/dnn\_call\_summary.csv} & the same for one call per multiply\\
\file{results/ibex/zeroskip\_summary.csv} & the zero-skipping / raw-pixel ablation, with full-test accuracy\\
\file{results/ibex/*\_images.csv} & one row per picture (answer, hash, cycles)\\
\file{results/ibex/profile.csv} & time per category for one traced picture\\
\bottomrule
\end{tabularx}
\end{table}

\section{How to rerun everything}

\begin{lstlisting}[language=bash]
# Phase 1: tests and GPU experiments
python3 -m pytest -q tests                      # 35 tests: PyTorch = Python = C
python3 experiments/run.py --dataset mnist   --stage all
python3 experiments/run.py --dataset fashion --stage all --epochs 30 --ft-epochs 5

# Phase 2: Ibex (one-time setup, no sudo), then all stages
hw/setup.sh
python3 hw/run_ibex.py --stage all --images 100 --jobs 48     # ~70 minutes on 48 cores
python3 hw/profile_ibex.py                                    # one traced picture per multiplier

# Slides and this document
python3 experiments/build_site.py               # index.html
cd doc && python3 make_figures.py && pdflatex roba_explained && pdflatex roba_explained
\end{lstlisting}
